\chapter{Background and Fundamental Concepts}


\section{Software Quality Models (ISO/IEC 25010)}

What is ISO/IEC 25010? According to Helpware.tech, ISO/IEC 25010 is a quality model that determines which characteristics to consider when evaluating the quality of a piece of software, whether its still in it's development phase, or even after it has already been released. 

This model consists of 8 different categories defining the quality of a software, from usability and reliability, to security and maintainability among others. These are the main categories this survey is going analyze and take into consideration when overviewing the AI tools.\cite{Quality}

\section{Technical Debt}

The term \textit{technical debt}, in software engineering, describes the cost that comes over periods of time from a suboptimal design or bad implementations that were made during the development phase. Such decision are usually for faster delivery of the product or for a less intense development effort, but usually results in a higher need of maintenance over time.

This phenomenon can happen for various reasons, such as missing documentation and comments, lacking testing, or a poorly structured code. Even though its not necessarily bad, an out of hand technical debt can affect the reliability and maintainability of a piece of software.

In our topic, AI-assisted development, technical debt becomes even more important. Due to automatically generated code, which prioritizes short-term functionality over long-term, if no proper human review is conducted the technical debt can accumulate, and so making future debugging more complex.

Because of this, technical debt is a concept that should be taken in consideration when long-term of AI in software engineering is analyzed.

\section{AI in Software Engineering: Definitions}

In software engineering, the term \textit{artificial intelligence} refers to computing systems designed to support tasks that are typically needed by humans. Some examples include patter recognition, decision-making and problem-solving. Unlike rule-based approaches, AI-based systems rely on data-based methods to adapt to multiple contexts.

Modern AI tools in software development are, more often than not, based on \textit{machine learning} techniques. Machine learning (or ML for short) models are trained on large datasets, often containing source code and the related documentation, in order to learn syntactic and semantic patterns. These models allow for tasks such as code completion, analysis, and automated recommendations.

In this survey the expression \textit{AI-based development tools} is used to describe tools that use AI techniques to help programmers during the development cycle. These tools are not made to operate autonomously, but to support programmers by giving suggestions, analyses, or automated actions.

This leads to AI being adopted in software engineering within a \textit{human-in-the-loop} paradigm, in which developers maintain control over the final decisions. This means that the effectiveness of these tools depends, not only on the quality of the models itself, but also on how the feedback they give back is interpreted and applied in practice.

\section{Ethical \& Legal Aspects in AI-Generated Code}

The use of AI in software engineering introduces ethical and legal challenges that is important they are considered along with the benefits. Particularly, the use of AI code generators may raise concerns regarding code ownership, responsibility as well as data privacy.

The most important legal issue has to do with copyright and intellectual property. Most AI models are trained on publicly available source code, often under different types of licenses. As a result, people wonder about the ownership and the legality of AI produced code.

Accountability is another important aspect to consider. When an AI tool applies directly changes to the source code, it's unclear who's responsible for bugs, vulnerabilities, or failures. Despite that, the final responsibility usually remains with the human developer, who's job is also to review and validate what an AI tool might've suggested.

Because of this, ethical and legal considerations play an important role in introducing AI tools. Reviewing these tools requires more than just a technical analysis, awareness of the risks they bring and their implications in professional development environments is also necessary.

